\chapter{ВЫБОР НАПРАВЛЕНИЯ ИССЛЕДОВАНИЙ И ОБЩАЯ МЕТОДИКА}
\label{cha:method}

\section{Обоснование направления исследований}

Разрабатываемый программный комплекс предназначен для моделирования и симуляции
робототехнических систем. Анализ предметной области и существующих решений
позволил выделить функциональные блоки, совокупность которых обеспечивает полный
цикл подготовки и проведения виртуального эксперимента:

\begin{enumerate}
  \item \textbf{Импорт и визуализация моделей.} Модели роботов в робототехнике
        традиционно описываются форматом URDF (Unified Robot Description Format),
        а сцены и объекты окружения --- форматами SDF (Simulation Description
        Format) и STL (стереолитография). Комплекс должен загружать эти форматы и
        отображать их в трёхмерной сцене в реальном времени.
  \item \textbf{Автоматическое построение кинематики.} Конструкторская
        документация на изделия чаще всего поставляется в нейтральном формате
        STEP (ISO~10303), который описывает геометрию сборки, но не содержит
        сведений о подвижных соединениях. Требуется алгоритм, восстанавливающий
        кинематическую схему (дерево звеньев и суставов) из геометрии сборки.
  \item \textbf{Физическая симуляция.} Ядро комплекса --- решатель динамики
        твёрдых тел с обнаружением столкновений и обработкой контактных
        ограничений. Для оптимизации трудозатрат принято решение в пользу
        единой кросс-платформенной реализации: один и тот же код исполняется
        как на центральном процессоре (CPU), так и на графических ускорителях
        (GPU) разных производителей, что исключает дублирование
        алгоритмической логики и расхождение версий.
  \item \textbf{Графический интерфейс пользователя.} Прикладная оболочка,
        объединяющая перечисленные подсистемы в единый рабочий процесс.
\end{enumerate}

\section{Сравнительный обзор существующих решений}

Задача физического моделирования твёрдых тел в робототехнике изучается с
1990-х~годов; основополагающие работы Д.~Барафа и А.~Уиткина~\cite{baraff1997}
заложили теоретическую базу контактной динамики. На сегодняшний день существует
ряд широко используемых систем~\cite{todorov2012mujoco,coumans2021bullet,nvidia_physx,smith2005ode}
(таблица~\ref{tab:engines}).

\begin{table}[ht]
  \centering
  \caption{Сравнение систем робототехнической симуляции}
  \label{tab:engines}
  \small
  \begin{tabular}{p{0.20\textwidth} p{0.30\textwidth} p{0.18\textwidth} p{0.18\textwidth}}
    \toprule
    Система & Физическое ядро & GPU-ускорение & Лицензия \\
    \midrule
    Gazebo / Ignition & ODE, Bullet, DART & Нет & Apache 2.0 \\
    MuJoCo            & Собственное (мягкие контакты) & Частично (MJX) & Apache 2.0 \\
    NVIDIA Isaac Sim  & PhysX 5 & Да & Проприетарная \\
    CoppeliaSim       & ODE, Bullet, Vortex, Newton & Нет & Проприетарная/\allowbreak GPL \\
    PyBullet          & Bullet & Нет & Zlib \\
    \bottomrule
  \end{tabular}
\end{table}

Перечисленные решения, как правило, не объединяют в одном продукте удобную среду
визуального моделирования, высокоточную симуляцию с массовым GPU-ускорением и
средства автоматической подготовки моделей из конструкторской документации.
Кроме того, ряд систем распространяется на условиях лицензий, ограничивающих
применение в государственном и коммерческом секторах. Это подтверждает
целесообразность разработки собственного комплекса.

\section{Методы решения задач и архитектура комплекса}

Прототип построен по модульному принципу. Подсистемы разделены по
ответственности и взаимодействуют через явно определённые программные
интерфейсы, что позволяет разрабатывать и тестировать их независимо
(рисунок~\ref{fig:arch}).

\begin{itemize}
  \item \texttt{modeuler} --- модуль загрузки и визуализации (C++); отрисовка
        выполняется средствами VulkanSceneGraph~\cite{vsg} поверх графического API Vulkan~\cite{vulkan};
  \item \texttt{kinechain} --- инструмент построения кинематической схемы из STEP
        (Python, геометрическое ядро OpenCASCADE~\cite{opencascade});
  \item \texttt{dyphur} --- физическое ядро (C++20) с переносимым вычислительным
        слоем поверх SYCL/AdaptiveCpp: единый исходный код компилируется для
        бэкендов CUDA (NVIDIA), HIP (AMD ROCm) и OpenMP (CPU);
  \item \texttt{modeuler-win} --- графический интерфейс пользователя (C\#, WPF).
\end{itemize}

\screenshot{0.7}{Архитектура прототипа программного комплекса}{Снимок экрана:
схема взаимодействия подсистем (modeuler, kinechain, dyphur, modeuler-win) ---
блок-схема компонентов и потоков данных между ними.}{arch}

\section{Общая методика проведения работ}

Разработка велась по итеративной методике с обязательным автоматизированным
тестированием. Для каждого модуля:
\begin{enumerate}
  \item формулировались функциональные требования и критерии приёмки;
  \item реализовывался алгоритм с покрытием модульными тестами;
  \item проводилось предварительное тестирование на искусственных
        (синтетических) данных, для которых известен ожидаемый результат;
  \item фиксировались эталонные показатели (корректность, детерминизм,
        производительность), используемые как опорные значения на последующих
        этапах.
\end{enumerate}

Воспроизводимость результатов обеспечивается контейнеризацией среды сборки и
фиксацией версий зависимостей. Особое внимание уделено \emph{детерминизму}
физического ядра: при фиксированных сборке и аппаратной платформе результаты
симуляции должны быть побитово воспроизводимыми, что является необходимым
условием для верификации алгоритмов управления и отладки.
